\chapter{研究背景與動機}
\label{ch:background}

% ==================================================
% 1.1 製程微縮與功耗
% ==================================================

\section{製程微縮與功耗}
\label{sec:power_budget}

隨著半導體製程技術持續微縮，
單位面積內可整合的電晶體數量不斷增加，
使積體電路得以實現更複雜的功能與更高的運算能力。
然而，電晶體密度的提升並不代表所有電路資源
都能在既定功率與散熱限制下同時運作。
傳統的 Dennard Scaling（定電場比例縮放）理論指出，
若電晶體尺寸與操作電壓能依適當比例同步縮小，
便有機會在提高整合密度的同時維持近似固定的功率密度。
然而，隨著製程技術進步，
供應電壓與臨界電壓的縮放逐漸受到元件效能、
漏電流及可靠度等因素限制，
使理想比例縮放難以持續。
當電晶體密度持續增加，
而功率密度無法得到有效控制時，
晶片便可能無法同時啟用所有運算資源，
形成暗矽（Dark Silicon）問題。
部分電路因此必須停用、降低操作頻率
或採取其他功率管理措施，
以符合系統的功率與散熱預算
\cite{dreslinski2010}。

在上述限制下，
供應電壓縮放（Supply-Voltage Scaling）
成為改善 CMOS 電路能量效率的重要方法。
CMOS 數位電路在邏輯狀態切換時，
需要對內部節點及負載電容進行充放電，
而供應電壓直接影響每次切換所消耗的能量。
因此，適當降低供應電壓能夠減少電路的切換能量，
並有機會降低系統的整體功耗。
然而，電壓降低也會減少電晶體的有效驅動電流，
使邏輯閘傳播延遲增加，
進而限制電路的最高操作頻率。
特別是在供應電壓接近電晶體臨界電壓時，
驅動能力對電壓變化更加敏感。
因此，近臨界運算（Near-Threshold Computing）
嘗試透過降低操作電壓，
以部分效能損失換取更低的切換能量，
但其實際效益仍取決於工作負載、
執行時間及漏電能量等因素
\cite{dreslinski2010,chandrakasan1992}。

不同應用系統對效能與能量的需求並不相同，
因此適合的供應電壓也有所差異。
例如，週期性運作的感測節點與生醫監測系統
通常具有相對寬鬆的吞吐量需求，
但受到電池容量與能量預算限制，
可在滿足任務完成期限的前提下，
利用較低的操作電壓減少切換能量。
相較之下，
視訊處理與其他即時運算系統
必須維持一定的資料處理速率，
但其工作負載可能隨時間變化。
若電路始終以最高效能需求所對應的供應電壓運作，
便可能在低負載期間消耗不必要的能量。
因此，設計者可以利用工作負載所提供的時序餘裕
調整供應電壓，
在滿足功能與效能要求的條件下改善能量效率
\cite{chandrakasan1992,udvs2010}。

值得注意的是，
製程微縮（Technology Scaling）
與供應電壓縮放並非相同的概念。
製程微縮涉及元件尺寸、結構及寄生效應的改變，
而供應電壓縮放則著重於操作電壓
對電路功耗、延遲與穩健性的影響。
Oldiges 等人針對 10/7-nm 技術的研究指出，
製程世代的功率與效能表現
受到元件結構及寄生效應等因素共同影響，
因此不能僅依製程節點名稱
判斷電路的節能幅度
\cite{oldiges2020}。
同樣地，在評估供應電壓縮放的效益時，
也必須區分電壓調整、
電路架構及製程技術所帶來的影響。
綜合而言，超低功耗 CMOS 設計的核心問題，
並非單純追求最低供應電壓，
而是如何在功能正確、時序與可靠度限制下，
選擇能夠有效降低任務能量的操作條件。


% ==================================================
% 1.2 CMOS 功率與能量基礎
% ==================================================

\section{CMOS 功率模型}
\label{sec:cmos_power_energy}

CMOS 數位電路的功率消耗主要可分為
動態切換功率（Dynamic Switching Power）、
短路功率（Short-Circuit Power）
與漏電功率（Leakage Power）。
其中，動態切換功率來自邏輯節點的充放電，
是供應電壓縮放能夠有效降低切換能量的主要原因。
然而，功率代表單位時間內的能量消耗，
而能量則取決於功率及其持續時間。
因此，在評估低電壓電路的節能效益時，
除了分析瞬時或平均功率，
也必須考慮完成指定運算所需的總能量。

當 CMOS 邏輯閘的輸出由低電位轉換至高電位時，
PMOS 電晶體提供充電路徑，
使負載電容 $C_L$
由零電位充電至供應電壓 $\VDD$。
在理想電壓源與固定電容的假設下，
電源提供的能量為 $C_L\VDD^2$，
其中一半以
$\frac{1}{2}C_L\VDD^2$
的形式儲存在負載電容，
另一半則於充電過程中耗散。
當輸出由高電位轉換至低電位時，
NMOS 電晶體提供放電路徑，
原本儲存在電容中的能量
也會透過放電路徑耗散。
因此，對一次完整的充放電循環而言，
由電源提供並最終耗散的能量為
$C_L\VDD^2$。

若以 $\alpha$ 表示每個時脈週期內的平均充電次數，
$f$ 表示時脈頻率，
則 CMOS 電路的平均切換功率可表示為

\begin{equation}
P_{\mathrm{sw}}
= \alpha C_L\VDD^2 f.
\label{eq:power}
\end{equation}

其中，$\alpha$ 為切換活動因子
（Switching Activity Factor），
用以描述節點的平均切換行為；
$C_L$ 為等效負載電容。
對於包含多個邏輯節點的實際電路，
可將各節點的切換功率加總，
或使用等效切換電容進行分析。
需要注意的是，
部分文獻將活動因子定義為
上升與下降轉換的總次數，
因此公式中可能出現
$\frac{1}{2}$ 的係數。
不同表示方式並不代表物理結果不同，
而是活動因子的定義有所差異
\cite{chandrakasan1992,gonzalez1997}。

由式~\eqref{eq:power} 可知，
當負載電容、切換活動因子與操作頻率固定時，
切換功率與供應電壓的平方成正比。
因此，若將供應電壓降低為原來的一半，
在上述條件不變的理想情況下，
切換功率可降低至原來的四分之一。
這種平方關係使供應電壓成為 CMOS
低功耗設計中特別重要的控制參數。
然而，實際降壓也會影響
電晶體的驅動能力與操作速度，
因此不能假設電路在降低電壓後
仍能維持原有的最高時脈頻率。

除了功率之外，
完成一次指定運算所需的能量
也是評估低功耗設計的重要指標。
假設某項工作需要 $N$ 個時脈週期，
且在執行期間具有固定的等效負載電容
與平均切換活動因子，
則總執行時間為
$T_{\mathrm{task}}=N/f$。
將切換功率乘以執行時間，
可以得到該工作的理想切換能量：

\begin{equation}
\begin{aligned}
E_{\mathrm{sw,task}}
&=P_{\mathrm{sw}}T_{\mathrm{task}}\\
&=N\alpha C_L\VDD^2.
\end{aligned}
\label{eq:switching_energy}
\end{equation}

式~\eqref{eq:switching_energy} 顯示，
在工作所需週期數與切換行為不變的條件下，
時脈頻率並不直接決定
完成固定任務的理想切換能量。
降低時脈頻率雖然能減少
單位時間內發生的切換次數，
使平均切換功率下降，
但也會延長任務的執行時間，
因此不會單獨降低
該任務的理想切換能量。
相較之下，
降低供應電壓能直接減少
每次電容充放電所需的能量。
然而，實際工作負載的切換活動、
執行週期數與電路狀態
可能隨操作條件改變，
因此上述結果應視為
固定工作量下的理想分析。

實際 CMOS 電路的功率消耗
並不只有負載電容的充放電。
當輸入訊號切換時，
PMOS 與 NMOS 電晶體
可能在短暫時間內同時導通，
形成由電源至接地的電流路徑，
產生短路功率。
此外，即使電晶體處於非切換狀態，
仍可能因次臨界導通、
接面漏電與閘極相關漏電機制
而消耗靜態功率。
因此，CMOS 電路的總功率可表示為

\begin{equation}
P_{\mathrm{total}}
=P_{\mathrm{sw}}
+P_{\mathrm{sc}}
+\VDD I_{\mathrm{leak}}.
\label{eq:total_power}
\end{equation}

其中，$P_{\mathrm{sc}}$ 為短路功率，
$I_{\mathrm{leak}}$ 則代表
指定電路範圍內的總漏電流。
短路功率通常受到
輸入訊號轉換時間、輸出負載
及元件特性影響；
漏電功率則與供應電壓、
溫度、製程及電路狀態有關。
因此，降低供應電壓雖然有助於
減少切換功率，
卻不代表總功率中的各個成分
都會按照相同的比例下降。

若進一步考慮完成指定任務
所需的總能量，則可表示為

\begin{equation}
\begin{aligned}
E_{\mathrm{task}}
={}&E_{\mathrm{sw,task}}
+E_{\mathrm{sc,task}}\\
&+\int_0^{T_{\mathrm{task}}}
\VDD(t)I_{\mathrm{leak}}(t)\,dt.
\end{aligned}
\label{eq:taskenergy}
\end{equation}

此式說明，
任務能量除了包含切換與短路能量之外，
也必須計入整段執行期間
所累積的漏電能量。
若供應電壓與漏電流
在執行期間近似固定，
漏電能量可簡化為

\begin{equation*}
E_{\mathrm{leak,task}}
\approx
\VDD I_{\mathrm{leak}}T_{\mathrm{task}}.
\end{equation*}

然而，當供應電壓降低
使電路操作速度下降時，
任務執行時間可能明顯增加，
導致漏電能量在總能量中
占有更高的比例。
因此，即使較低的電壓
能夠有效降低切換能量，
也不一定能使完成指定任務
所需的總能量持續下降。

綜合上述分析，
供應電壓縮放的節能效益
必須同時考慮切換能量、
短路能量、漏電能量
及任務執行時間。
尤其在近臨界與次臨界操作區域，
降低電壓所帶來的切換能量改善，
可能與執行時間延長
所增加的漏電能量形成取捨。
因此，超低功耗 CMOS 設計的目標
不應只追求最低瞬時功率
或最低操作電壓，
而應在滿足功能與時序需求的條件下，
尋找適當的供應電壓操作點。
後續章節將進一步討論
不同操作區域的特性，
以及影響最低能量點的物理限制。

\section{操作區域與目標}
\label{sec:operating_regions}

供應電壓的降低會改變 MOSFET 的操作特性，
進而影響 CMOS 電路的切換速度、
能量消耗與功能穩健性。
依據供應電壓相對於電晶體臨界電壓
$V_{\mathrm{th}}$ 的大小，
低功耗 CMOS 電路通常可概略區分為
超臨界操作（Superthreshold Operation）、
近臨界操作（Near-Threshold Operation）
與次臨界操作（Subthreshold Operation）
\cite{dreslinski2010,calhoun2005,markovic2010}。

超臨界操作通常採用明顯高於臨界電壓的
供應電壓，使電晶體具有較大的有效驅動電流，
因此能夠支援較高的操作頻率。
然而，較高的供應電壓也會增加
每次電容充放電所消耗的能量。
因此，在效能需求允許的情況下，
適當降低供應電壓可有效改善切換能量。

近臨界操作則將供應電壓降低至
與臨界電壓相近的範圍。
相較於傳統較高電壓的操作方式，
近臨界運算能顯著降低切換能量，
但電晶體的驅動能力也會快速下降，
使電路延遲增加。
此外，當供應電壓接近臨界電壓時，
製程、電壓與溫度變異對電路效能的影響
可能更加明顯。
因此，近臨界運算的設計重點，
是在能量改善、操作速度與功能穩健性之間
取得適當的平衡。

次臨界操作使用低於臨界電壓的
供應電壓，使電晶體主要依靠
弱反轉（Weak Inversion）電流完成切換。
雖然此時電晶體仍能導通，
但可用的驅動電流相當有限，
使邏輯閘傳播延遲大幅增加。
因此，次臨界操作通常較適合
吞吐量需求較低、允許較長執行時間的應用。
然而，過長的執行時間也可能使漏電能量
在總任務能量中占據較大的比例，
因此最低操作電壓不一定對應最低任務能量。

需要注意的是，
上述操作區域並不存在適用於所有製程的
固定供應電壓分界。
實際操作特性會受到臨界電壓定義、
元件結構、溫度與偏壓條件影響。
因此，相同的供應電壓在不同製程中，
可能對應不同的操作區域。
此外，數位電路中的電晶體
會隨邏輯狀態而具有不同的瞬時偏壓，
因此操作區域的分類應視為
描述電路整體操作條件的近似方式，
而非所有元件在任意時刻皆處於相同反轉狀態。

除了操作區域的差異，
低功耗設計也必須明確定義最佳化目標。
若僅以每次任務所消耗的能量為目標，
設計者便會尋找最低能量點
（Minimum-Energy Point, MEP）。
然而，當應用同時重視運算速度時，
也可以採用能量延遲乘積
（Energy-Delay Product, EDP）
作為評估指標：

\begin{equation}
EDP = E_{\mathrm{task}} T_{\mathrm{task}}.
\label{eq:energy_delay_product}
\end{equation}

其中，$E_{\mathrm{task}}$ 為完成指定任務
所需的總能量，
$T_{\mathrm{task}}$ 為任務執行時間。
由於 EDP 同時考慮能量與延遲，
其最佳操作電壓未必與最低能量點相同。
Gonzalez 等人的研究著重於
能量延遲最佳化，
而 Calhoun 等人則分析
低電壓電路的最低操作能量
\cite{gonzalez1997,calhoun2005}。
因此，比較不同研究結果時，
必須先確認其最佳化目標與評估條件。

對於具有明確完成期限的實際應用，
更合理的設計方式是在滿足功能正確性、
時序與可靠度要求的前提下，
最小化完成指定任務所需的總能量。
此時，最佳供應電壓除了受到
電晶體切換特性影響，
也可能受到記憶體最低操作電壓、
系統吞吐量及電源轉換損失限制。
因此，超低功耗 CMOS 設計
並非單純選擇最低供應電壓，
而是需要依據工作負載與系統限制，
尋找具有實際可行性的最佳操作點。


\section{文獻與研究架構}
\label{sec:literature_framework}

本報告以七篇 CMOS 低功耗設計相關研究論文
作為主要文獻依據，
從電晶體的物理特性、電路操作限制，
到架構與系統層級的最佳化方法，
探討供應電壓縮放的節能效益與實作挑戰。

在基礎理論方面，
Chandrakasan 等人探討 CMOS 數位電路的
低功耗設計方法，
說明供應電壓與電路架構對功率消耗的影響
\cite{chandrakasan1992}。
Gonzalez 等人進一步研究
供應電壓與臨界電壓的共同最佳化，
分析能量延遲乘積與操作條件之間的關係
\cite{gonzalez1997}。
這些研究建立了電壓縮放與
能量效率分析的理論基礎。

在近臨界與次臨界運算方面，
Calhoun 等人透過低電壓 FIR 濾波器研究，
分析最低能量點及其對設計條件的敏感性
\cite{calhoun2005}。
Dreslinski 等人則從處理器與系統角度，
整理近臨界運算的節能潛力，
以及製程變異、記憶體與效能下降
所帶來的設計挑戰
\cite{dreslinski2010}。
Marković 等人的研究進一步探討
超低電壓數位電路的設計方法，
包含元件、邏輯與架構層級的能量最佳化策略
\cite{markovic2010}。

在系統整合與製程技術方面，
超寬範圍動態電壓縮放
（Ultra-Dynamic Voltage Scaling, UDVS）
相關研究討論低電壓 SRAM、
電路輔助技術與電源管理設計，
說明單純降低邏輯電路的操作電壓
不足以保證整體系統具有最佳能量效率
\cite{udvs2010}。
此外，Oldiges 等人針對
10/7-nm 技術的研究，
說明元件結構、寄生效應與製程特性
對功率及效能縮放的影響
\cite{oldiges2020}。

根據上述研究，
本報告首先分析供應電壓降低後
所面臨的驅動電流下降、
延遲增加、PVT 變異與 SRAM 穩定性問題。
接著依照元件、電路及架構與系統
三個設計層級，
探討改善低電壓操作能力與能量效率的方法。
在此基礎上，
進一步比較實際晶片的量測結果，
並整理低電壓電路的實作與驗證流程。
最後歸納適用於不同工作負載的設計準則，
並討論先進元件與自適應控制技術的發展方向。

透過由物理限制到多層次設計方法的分析，
本報告希望說明：
供應電壓縮放的價值不僅在於降低切換能量，
更在於如何結合電路特性、
系統需求與實作限制，
選擇兼顧能量效率、效能與可靠度的操作條件。
